\documentclass[12pt]{article}
\usepackage[utf8]{inputenc}
\usepackage[T1]{fontenc}
\usepackage{hyperref}
\title{So Language Resources for K-12 Curriculum Development}
\author{Compiled by Claude Sonnet 5.5 (claude-sonnet-5-5) for Andr\'as Kornai}
\date{}
\begin{document}
\maketitle

\section{Introduction}
So (ISO 639-3 sss; Glottocode sooo1254) is a Katuic language of the Austroasiatic family, spoken in Laos and Thailand. The people are known locally as So, Bru or Katang. The name ``So'' is ambiguous. My search surfaced material on Thavung (also called So (Thavung)), which is a separate Vietic language of the same area. I do not list it as a resource for sss. Very little specific to this language was found, and nothing on schooling.

\section{Primary Educational Resources}
I found no primers, textbooks or school programmes for So.

\section{Linguistic and Reference Materials}
\begin{itemize}
\item ASJP language page for So (Laos), a word list with a speaker estimate: \url{https://asjp.clld.org/languages/SO_LAOS}
\item Migliazza, paper on ethnolinguistic groups of the area; it estimates about 35,000 So in Thailand and about 100,000 in central Laos: \url{https://www.ethnorema.it/wp-content/uploads/1-03-Migliazza.pdf}
\end{itemize}

\section{Religious and Community Resources}
None found.

\section{Digital and Audio Resources}
None found for So proper. A Thavung lexical resource also appeared (\url{https://cocoon.huma-num.fr/exist/crdo/lexvo.xql?code=thm}), but it concerns a different language and is mentioned only as a warning against confusion.

\section{Curriculum Development Context}
Sources describe language shift toward Thai and the northeastern dialect in Thailand. Speaker figures differ greatly between sources (about 35,000 in Thailand, 100,000 to 233,000 in Laos). No mother-tongue education work for So was found. A curriculum would start from scratch, with community and linguist partners.

\end{document}